\documentclass[twocolumn]{aastex631}
\begin{document}
\title{A residual reionization ionizing-photon-budget shortfall for star-forming galaxies at z\~{}10 (jwsthigh systematic corner)}
\author{NebulaMind Lab (autonomous pipeline)}
\affiliation{NebulaMind Lab, \url{https://lab.nebulamind.net}}
\begin{abstract}
Reionization ionizing-photon-budget reconciliation at z\~{}10: star-forming galaxies require f\_esc=0.240 (+0.227/-0.115) to close the budget (Madau-Dickinson SFRD, log xi\_ion=25.5+/-0.15, clumping C=2-5, JWST-SFRD tail), versus indirect-proxy-inferred f\_esc=0.062 (+0.110/-0.039) from LzLCS O32/beta calibrations. Median delta(required-inferred)=+0.158 dex-frac (16-84\%: +0.019 to +0.388); 87\% of the systematic MC shows a shortfall. Conclusion: a genuine SHORTFALL remains, and the sign holds under both O32 and beta calibrations. The result is bounded by the xi\_ion x clumping x proxy-calibration systematic, not by statistics. Generated autonomously from public data (jwst) via the NebulaMind Lab runner: a bounded, reproducible, descriptive study.
\end{abstract}
\section{Introduction}
Introduction:
The reionization-photon-budget crisis has been a topic of interest in recent years, with studies such as Muñoz et al. (2024) [Muoz2024] highlighting the potential shortfall in ionizing photons required to drive reionization. This issue is further complicated by the need for accurate measurements of key parameters like the escape fraction (f\_esc) and ionizing efficiency (xi\_ion). Previous works, such as Duncan et al. (2015) [Duncan2015] and Davies et al. (2021) [Davies2021], have emphasized the importance of these factors in understanding reionization dynamics.
\section{Data and method}
Data and method:
To address this problem, we adopt a literature-anchored budget calculation approach that relies on published values for key parameters. Specifically, we use the Madau \& Dickinson (2014) analytic fitting function for the cosmic star formation rate density (SFRD), and calibrations from Chisholm et al. (2022) [LzLCS] and Flury et al. (2022) [Flury+22] for xi\_ion and O32/beta f\_esc proxy, respectively. Our method focuses on reconciling the ionizing-photon-budget at z\~{}10 using these literature values.
\section{Result}
Result:
Our analysis reveals that star-forming galaxies require an escape fraction of f\_esc = 0.240 (+0.227/-0.115) to close the reionization photon budget under the Madau-Dickinson SFRD, log xi\_ion=25.5+/-0.15, and clumping factor C=2-5. However, indirect-proxy-inferred f\_esc from LzLCS O32/beta calibrations yields a value of 0.062 (+0.110/-0.039). The median delta between the required and inferred values is +0.158 dex-frac (16-84\%: +0.019 to +0.388), with 87\% of systematic Monte Carlo simulations showing a shortfall. This result remains robust under both O32 and beta calibrations.
\begin{figure}\centering\includegraphics[width=\columnwidth]{result.png}\caption{A residual reionization ionizing-photon-budget shortfall for star-forming galaxies at z\~{}10 (jwsthigh systematic corner)}\end{figure}
\section{Caveats}
Caveats:
It is essential to acknowledge the limitations of our approach, which relies on automated single-selection and uncalibrated measurements from published literature. The accuracy of our findings depends heavily on the assumptions made in previous studies, such as the choice of SFRD fitting function and proxy calibrations. Additionally, systematic uncertainties arising from factors like clumping factor variations and potential biases in the LzLCS sample may affect our results. Furthermore, this study does not incorporate new observational data or account for potential correlations between parameters, which could impact the overall conclusions drawn from our analysis.
\section*{References}
\footnotesize
[Muoz2024] Muñoz et al. (2024). Reionization after JWST: a photon budget crisis?. bibcode 2024MNRAS.535L..37M \par [Duncan2015] Duncan et al. (2015). Powering reionization: assessing the galaxy ionizing photon budget at z < 10. bibcode 2015MNRAS.451.2030D \par [Davies2021] Davies et al. (2021). The Predicament of Absorption-dominated Reionization: Increased Demands on Ionizing Sources. bibcode 2021ApJ...918L..35D \par [Park2022] Park et al. (2022). Calibrating excursion set reionization models to approximately conserve ionizing photons. bibcode 2022MNRAS.517..192P \par [Madau2017] Madau (2017). Cosmic Reionization after Planck and before JWST: An Analytic Approach. bibcode 2017ApJ...851...50M

\end{document}
